% Chapter 3, Topic _Linear Algebra_ Jim Hefferon
%  http://joshua.smcvt.edu/linearalgebra
%  2001-Jun-12
\topic{Matriks Ortonormal}
\index{matriks ortonormal|(}
\index{matriks!ortonormal|(}
Dalam \emph{The Elements},\index{Euklides} Euklides menganggap dua bangun
sama jika keduanya memiliki ukuran dan bentuk yang sama.
Jadi, meskipun segitiga-segitiga di bawah ini tidak sama karena bukan himpunan
titik yang sama, untuk keperluan Euklides keduanya pada dasarnya tidak dapat
dibedakan. Kita dapat membayangkan mengangkat bidangnya, menggeser dan sedikit
memutarnya tanpa membengkokkan atau meregangkannya, lalu meletakkannya kembali
sehingga bangun pertama berimpit dengan bangun kedua.
(Euklides tidak pernah menyatakan asas ini secara eksplisit, tetapi sering
menggunakannya \cite{Casey}.)
\begin{center}
  \includegraphics{map/mp/ch3.61}
\end{center}
Dalam istilah modern, ``mengangkat bidang~\ldots'' berarti mengambil pemetaan
dari bidang ke dirinya sendiri.
Euklides hanya meninjau transformasi yang dapat menggeser atau memutar bidang,
tetapi tidak membengkokkan atau meregangkannya.
Karena itu, pemetaan $\map{f}{\Re^2}{\Re^2}$ disebut
\definend{pemetaan yang mempertahankan jarak}%
\index{pemetaan yang mempertahankan jarak}%
\index{pemetaan!mempertahankan jarak}%
\index{fungsi!mempertahankan jarak},
\definend{gerak kaku}\index{gerak kaku}, atau
\definend{isometri}\index{isometri}, jika untuk semua titik
$P_1,P_2\in\Re^2$, jarak dari $f(P_1)$ ke $f(P_2)$ sama dengan jarak dari
$P_1$ ke $P_2$.
Kita juga mendefinisikan \definend{bangun} bidang\index{bangun bidang} sebagai
suatu himpunan titik pada bidang dan mengatakan bahwa dua bangun
\definend{kongruen}\index{bangun bidang!kongruen}%
\index{kongruensi!bangun bidang} jika terdapat pemetaan yang mempertahankan
jarak dari bidang ke dirinya sendiri yang membawa satu bangun ke bangun lainnya.

Berbagai pernyataan dalam geometri Euklides mudah mengikuti definisi-definisi
ini. Beberapa di antaranya adalah: (i)~kekolinearan invarian terhadap sebarang
pemetaan yang mempertahankan jarak (artinya, jika $P_1$, $P_2$, dan $P_3$
kolinear, maka $f(P_1)$, $f(P_2)$, dan $f(P_3)$ juga kolinear),
(ii)~keterantaraan invarian terhadap sebarang pemetaan yang mempertahankan
jarak (jika $P_2$ berada di antara $P_1$ dan $P_3$, maka $f(P_2)$ berada di
antara $f(P_1)$ dan $f(P_3)$), (iii)~sifat sebagai segitiga invarian terhadap
sebarang pemetaan yang mempertahankan jarak (jika suatu bangun merupakan
segitiga, maka bayangannya juga merupakan segitiga), dan (iv)~sifat sebagai
lingkaran invarian terhadap sebarang pemetaan yang mempertahankan jarak.
Pada tahun 1872, F.~Klein\index{Klein, F.} menyarankan agar geometri Euklides
didefinisikan sebagai kajian sifat-sifat yang invarian terhadap
pemetaan-pemetaan tersebut.
(Gagasan ini merupakan bagian dari Program Erlangen Klein
\index{Program Erlangen}, yang mengusulkan asas pengorganisasian bahwa setiap
jenis geometri\Dash Euklides, proyektif, dan sebagainya\Dash dapat
dideskripsikan sebagai kajian sifat-sifat yang invarian terhadap suatu grup
transformasi. Kata `grup' di sini berarti lebih dari sekadar `kumpulan', tetapi
hal itu berada di luar cakupan kita.)

Kita dapat menggunakan aljabar linear untuk mencirikan pemetaan-pemetaan bidang
yang mempertahankan jarak.

Sebagai permulaan, perhatikan bahwa terdapat transformasi bidang yang
mempertahankan jarak tetapi tidak linear.
Contoh yang jelas adalah \emph{translasi} berikut.\index{translasi}
\begin{equation*}
  \colvec{x \\ y}
  \quad\mapsto\quad
  \colvec{x \\ y}+\colvec[r]{1 \\ 0}=\colvec{x+1 \\ y}
\end{equation*}
Namun, ternyata contoh ini merupakan satu-satunya pengecualian, dalam arti
bahwa jika $f$ mempertahankan jarak dan membawa $\zero$ ke $\vec{v}_0$, maka
pemetaan $\vec{v}\mapsto f(\vec{v})-\vec{v}_0$ bersifat linear.
Hal ini langsung mengikuti pernyataan berikut:~pemetaan $t$ yang
mempertahankan jarak dan membawa $\zero$ ke dirinya sendiri bersifat linear.
Untuk membuktikan pernyataan ekuivalen ini, tinjau basis standar dan andaikan
\begin{equation*}
  t(\vec{e}_1)=\colvec{a \\ b}
  \qquad
  t(\vec{e}_2)=\colvec{c \\ d}
\end{equation*}
untuk suatu $a,b,c,d\in\Re$.
Untuk menunjukkan bahwa $t$ linear, kita dapat menunjukkan bahwa $t$ dapat
direpresentasikan oleh matriks; artinya, $t$ bertindak sebagai berikut untuk
semua $x,y\in\Re$.
\begin{equation*}
  \vec{v}=\colvec{x  \\  y}
  \mapsunder{t}
  \colvec{ax+cy  \\  bx+dy}
\tag*{\text{($*$)}}\end{equation*}
Ingat bahwa jika kita menetapkan tiga titik yang tidak kolinear, sebarang titik
dapat ditentukan dengan memberikan jaraknya dari ketiga titik tersebut.
Jadi, kita dapat menentukan sebarang titik $\vec{v}$ dalam domain berdasarkan
jaraknya dari $\zero$, $\vec{e}_1$, dan $\vec{e}_2$.
Demikian pula, kita dapat menentukan sebarang titik $t(\vec{v})$ dalam
kodomain berdasarkan jaraknya dari ketiga titik tetap $t(\zero)$,
$t(\vec{e}_1)$, dan $t(\vec{e}_2)$ (ketiganya tidak kolinear karena, seperti
disebutkan di atas, kekolinearan bersifat invarian, sedangkan $\zero$,
$\vec{e}_1$, dan $\vec{e}_2$ tidak kolinear).
Karena $t$ mempertahankan jarak, kita dapat menyatakan lebih banyak:~untuk
titik $\vec{v}$ pada bidang yang ditentukan oleh jarak $d_0$ dari $\zero$,
jarak $d_1$ dari $\vec{e}_1$, dan jarak $d_2$ dari $\vec{e}_2$, bayangannya
$t(\vec{v})$ harus merupakan titik tunggal dalam kodomain yang ditentukan oleh
jarak $d_0$ dari $t(\zero)$, $d_1$ dari $t(\vec{e}_1)$, dan $d_2$ dari
$t(\vec{e}_2)$.
Karena ketunggalan tersebut, cukup periksa bahwa tindakan dalam ($*$) berlaku
untuk kasus $d_0$, $d_1$, dan $d_2$:
\begin{equation*}
  \dist(\colvec{x \\ y},\zero) 
  =\dist(t(\colvec{x \\ y}),t(\zero))
  =\dist(\colvec{ax+cy \\ bx+dy},\zero)
\end{equation*}
(kita mengasumsikan bahwa $t$ memetakan $\zero$ ke dirinya sendiri)
\begin{equation*}
  \dist(\colvec{x \\ y},\vec{e}_1)
  =\dist(t(\colvec{x \\ y}),t(\vec{e}_1))
  =\dist(\colvec{ax+cy \\ bx+dy},\colvec{a \\ b})  
\end{equation*}
dan
\begin{equation*}
  \dist(\colvec{x \\ y},\vec{e}_2)
  =\dist(t(\colvec{x \\ y}),t(\vec{e}_2))
  =\dist(\colvec{ax+cy \\ bx+dy},\colvec{c \\ d})
\end{equation*}
Pemeriksaan-pemeriksaan tersebut menunjukkan bahwa ($*$) mendeskripsikan $t$;
semuanya bersifat rutin.

Jadi, sebarang $\map{f}{\Re^2}{\Re^2}$ yang mempertahankan jarak merupakan
pemetaan linear ditambah translasi,
$f(\vec{v})=t(\vec{v})+\vec{v}_0$, untuk suatu vektor konstan $\vec{v}_0$ dan
pemetaan linear $t$ yang mempertahankan jarak.
Karena itu, untuk memahami pemetaan yang mempertahankan jarak, yang tersisa
adalah memahami pemetaan linear yang mempertahankan jarak.

Tidak setiap pemetaan linear mempertahankan jarak.
Sebagai contoh, $\vec{v}\mapsto 2\vec{v}$ tidak mempertahankan jarak.

Namun, terdapat pencirian yang rapi:~transformasi linear $t$ pada bidang
mempertahankan jarak jika dan hanya jika
$\norm{t(\vec{e}_1)}=\norm{t(\vec{e}_2)}=1$ dan $t(\vec{e}_1)$ ortogonal
terhadap $t(\vec{e}_2)$.
Arah `hanya jika' dari pernyataan tersebut mudah\Dash karena $t$
mempertahankan jarak, $t$ harus mempertahankan panjang vektor; dan menurut
teorema Pythagoras, $t$ juga harus mempertahankan keortogonalan.
Untuk menunjukkan arah `jika', kita dapat memeriksa bahwa pemetaan tersebut
mempertahankan panjang vektor, sebab dengan demikian jarak antara setiap
$\vec{p}$ dan $\vec{q}$ dipertahankan:
$\norm{t(\vec{p}-\vec{q}\,)}=\norm{t(\vec{p})-t(\vec{q}\,)}
=\norm{\vec{p}-\vec{q}\,}$.
Untuk pemeriksaan tersebut, misalkan
\begin{equation*}
  \vec{v}=\colvec{x \\ y} 
  \quad 
  t(\vec{e}_1)=\colvec{a \\ b}
  \quad
  t(\vec{e}_2)=\colvec{c \\ d}
\end{equation*}
dan berdasarkan asumsi arah `jika', yaitu
$a^2+b^2=c^2+d^2=1$ dan $ac+bd=0$, kita memiliki
\begin{align*}
  \norm{t(\vec{v}\,)}^2
  &= (ax+cy)^2+(bx+dy)^2  \\
  &= a^2x^2+2acxy+c^2y^2+b^2x^2+2bdxy+d^2y^2 \\
  &= x^2(a^2+b^2)+y^2(c^2+d^2)+2xy(ac+bd)  \\
  &= x^2 + y^2 \\
  &= \norm{\vec{v}\,}^2
\end{align*}

Salah satu hal menarik dari pencirian ini adalah bahwa kita dapat dengan mudah
mengenali matriks yang merepresentasikan pemetaan semacam itu terhadap basis
standar:~kolom-kolomnya memiliki panjang satu dan saling ortogonal.
Matriks semacam ini disebut \definend{matriks ortonormal}
\index{matriks ortonormal}\index{matriks!ortonormal} (atau, secara kurang
formal, \definend{matriks ortogonal}\index{matriks ortogonal}%
\index{matriks!ortogonal}, karena istilah ini sering digunakan bukan hanya
untuk menyatakan bahwa kolom-kolomnya ortogonal, tetapi juga bahwa panjangnya
satu).

Kita dapat memanfaatkan pencirian ini untuk memahami tindakan geometris
pemetaan-pemetaan yang mempertahankan jarak.
Karena $\norm{t(\vec{v}\,)}=\norm{\vec{v}\,}$, pemetaan~$t$ membawa sebarang
$\vec{v}$ ke suatu titik pada lingkaran berpusat di titik asal dengan jari-jari
yang sama dengan panjang $\vec{v}$.
Khususnya, $\vec{e}_1$ dan $\vec{e}_2$ dipetakan ke lingkaran satuan.
Lebih lanjut, %because of the orthogonality restriction,
setelah kita menetapkan bahwa vektor satuan $\vec{e}_1$ dipetakan ke vektor
dengan komponen $a$ dan $b$, hanya terdapat dua tempat tujuan $\vec{e}_2$ jika
bayangannya harus tegak lurus terhadap bayangan vektor pertama:~$\vec{e}_2$
dapat dipetakan dengan mempertahankan posisinya seperempat lingkaran berlawanan arah
jarum jam dari $\vec{e}_1$
\begin{center}
  \begin{tabular}{@{}c@{}}\includegraphics{map/mp/ch3.62}\end{tabular}
  \qquad
  $\rep{t}{\stdbasis_2,\stdbasis_2}=
  \begin{mat}
    a  &-b  \\
    b  &a
  \end{mat}$ 
\end{center}
atau ke posisi seperempat lingkaran searah jarum jam.
\begin{center}
  \begin{tabular}{@{}c@{}}\includegraphics{map/mp/ch3.63}\end{tabular}
  \qquad
  $\rep{t}{\stdbasis_2,\stdbasis_2}=
  \begin{mat}
    a  &b  \\
    b  &-a
  \end{mat}$
\end{center}

Deskripsi geometris kedua kasus ini mudah.
Misalkan $\theta$ merupakan sudut berlawanan arah jarum jam antara sumbu-$x$
dan bayangan $\vec{e}_1$.
Matriks pertama di atas merepresentasikan, terhadap basis standar,
\definend{rotasi}\index{rotasi}\index{pemetaan linear!rotasi} bidang sebesar
$\theta$ radian.
\begin{center}
  \begin{tabular}{@{}c@{}}\includegraphics{map/mp/ch3.62}\end{tabular}
  \qquad
  $\colvec{x \\ y}
  \mapsunder{t}
  \colvec{x\cos\theta-y\sin\theta  \\ x\sin\theta+y\cos\theta}$ 
\end{center}
Matriks kedua di atas merepresentasikan
\definend{pencerminan}\index{pencerminan}%
\index{pemetaan linear!pencerminan} bidang terhadap garis yang membagi dua
sudut antara $\vec{e}_1$ dan $t(\vec{e}_1)$.
\begin{center}
  \begin{tabular}{@{}c@{}}\includegraphics{map/mp/ch3.64}\end{tabular}
  \qquad
  $\colvec{x \\ y}\mapsunder{t}
  \colvec{x\cos\theta+y\sin\theta  \\  x\sin\theta-y\cos\theta}$ 
\end{center}
(Gambar ini menunjukkan $\vec{e}_1$ dicerminkan ke atas menuju kuadran pertama
dan $\vec{e}_2$ dicerminkan ke bawah menuju kuadran keempat.)

Perhatikan:~dalam domain, sudut dari $\vec{e}_1$ ke $\vec{e}_2$ berjalan
berlawanan arah jarum jam, dan pada pemetaan pertama di atas, sudut dari
$t(\vec{e}_1)$ ke $t(\vec{e}_2)$ juga berjalan berlawanan arah jarum jam,
sehingga pemetaan tersebut mempertahankan orientasi sudut.
Namun, pemetaan kedua membalik orientasinya.
Pemetaan yang mempertahankan jarak disebut
\definend{isometri langsung}\index{isometri!langsung}%
\index{pemetaan!mempertahankan orientasi} jika mempertahankan orientasi, dan
\definend{isometri berlawanan}\index{isometri!berlawanan}%
\index{pemetaan!membalik orientasi} jika membalik orientasi.

Dengan demikian, kita telah mencirikan kajian Euklides tentang kongruensi.
Bagi bangun bidang, kajian tersebut meninjau sifat-sifat yang invarian terhadap
kombinasi (i)~rotasi yang diikuti translasi, atau (ii)~pencerminan yang diikuti
translasi (pencerminan yang diikuti translasi nontrivial sejajar dengan sumbu
pencerminan disebut
\definend{refleksi geser}\index{pencerminan!refleksi geser}).

Gagasan lain yang dijumpai dalam geometri elementer, selain kongruensi bangun,
adalah bahwa bangun-bangun disebut
\definend{sebangun}\index{segitiga sebangun}\index{segitiga!sebangun} jika
keduanya kongruen setelah perubahan skala.
Kedua segitiga di bawah ini sebangun karena segitiga kedua memiliki bentuk yang
sama dengan segitiga pertama, tetapi berukuran $3/2$ kali ukuran segitiga pertama.
\begin{center}
  \includegraphics{map/mp/ch3.65}
\end{center}
Dari pembahasan di atas, dua bangun sebangun jika terdapat matriks ortonormal
$T$ sedemikian sehingga titik-titik $\vec{q}$ pada satu bangun merupakan
bayangan titik-titik $\vec{p}$ pada bangun lainnya melalui
$\vec{q}=(kT)\vec{p}+\vec{p}_0$, untuk suatu bilangan riil tak nol $k$ dan
vektor konstan $\vec{p}_0$.

Meskipun gagasan-gagasan ini berasal dari Euklides, matematika tidak lekang
oleh waktu dan semuanya masih digunakan saat ini.
Salah satu penerapan pemetaan yang dikaji di atas terdapat dalam grafika
komputer. Sebagai contoh, kita dapat menganimasikan tampak atas kubus berikut
dengan menyusun bingkai-bingkai film saat kubus itu berputar; gerakan tersebut
merupakan gerak kaku.
\begin{center}
  \begin{tabular}{ccc}
    \includegraphics{map/mp/ch3.66}
    &\includegraphics{map/mp/ch3.67}
    &\includegraphics{map/mp/ch3.68}                       \\
    Bingkai 1      &Bingkai 2       &Bingkai 3
  \end{tabular}
\end{center}
Kita juga dapat membuat kubus tampak menjauh dengan menghasilkan
bingkai-bingkai film yang memperlihatkannya mengecil, sehingga memberikan
bangun-bangun yang sebangun.
\begin{center}
  \begin{tabular}{ccc}
    \includegraphics{map/mp/ch3.69}
    &\includegraphics{map/mp/ch3.70}
    &\includegraphics{map/mp/ch3.71}                       \\
    Bingkai 1:        &Bingkai 2:       &Bingkai 3:
  \end{tabular}
\end{center}
Grafika komputer memanfaatkan teknik-teknik aljabar linear dalam banyak cara
lain (lihat \nearbyexercise{exer:HomoCrds}).

% So the analysis above of distance-preserving maps
% is useful as well as interesting.
Buku indah yang menjelajahi sebagian bidang ini adalah \cite{Weyl}.
Pembahasan lebih lanjut tentang grup, baik grup transformasi maupun yang lain,
terdapat dalam berbagai buku Aljabar Modern, misalnya
\cite{BirkhoffMaclane}. Pembahasan lebih lanjut tentang Klein dan Program
Erlangen terdapat dalam \cite{Yaglom}.
 

\begin{exercises}
  \item 
    Tentukan apakah setiap matriks berikut ortonormal.
    \begin{exparts}
      \partsitem $\begin{mat}[r]
                    1/\sqrt{2} &-1/\sqrt{2}  \\
                   -1/\sqrt{2} &-1/\sqrt{2}
                  \end{mat}$
      \partsitem $\begin{mat}[r]
                    1/\sqrt{3} &-1/\sqrt{3}  \\
                   -1/\sqrt{3} &-1/\sqrt{3}
                  \end{mat}$
      \partsitem $\begin{mat}[r]
                    1/\sqrt{3} &-\sqrt{2}/\sqrt{3}  \\
                   -\sqrt{2}/\sqrt{3} &-1/\sqrt{3}
                  \end{mat}$
    \end{exparts}
    \begin{answer}
      \begin{exparts}
        \partsitem Ya.
        \partsitem Tidak, panjang kolom-kolomnya bukan satu.
        \partsitem Ya.
      \end{exparts}
    \end{answer}
  \item 
    Tuliskan rumus bagi setiap pemetaan yang mempertahankan jarak berikut.
    \begin{exparts}  
      \partsitem pemetaan yang merotasi sebesar $\pi/6$ radian, lalu
        mentranslasi sebesar $\vec{e}_2$
      \partsitem pemetaan yang mencerminkan terhadap garis $y=2x$
      \partsitem pemetaan yang mencerminkan terhadap $y=-2x$, lalu
        mentranslasi satu satuan ke kanan dan satu satuan ke atas
    \end{exparts}
    \begin{answer}
      Beberapa di antaranya tidak linear karena melibatkan translasi
      nontrivial.
      \begin{exparts}
        \partsitem 
          $\colvec{x \\ y}
           \mapsto
           \begin{mat}
             x\cdot\cos(\pi/6)-y\cdot\sin(\pi/6) \\
             x\cdot\sin(\pi/6)+y\cdot\cos(\pi/6) 
           \end{mat}
           +\colvec[r]{0 \\ 1}
           =\begin{mat}
             x\cdot(\sqrt{3}/2)-y\cdot(1/2)+0 \\
             x\cdot(1/2)+y\cdot(\sqrt{3}/2)+1
           \end{mat}$
        \partsitem Garis $y=2x$ membentuk sudut $\alpha=\arctan(2/1)$
          terhadap sumbu-$x$.
          Karena matriks pencerminan menggunakan sudut $2\alpha$, kita
          memperoleh $\cos(2\alpha)=-3/5$ dan $\sin(2\alpha)=4/5$.
          \begin{equation*}
          \colvec{x \\ y}
           \mapsto
           \begin{mat}
             -3x/5+4y/5 \\
             4x/5+3y/5
           \end{mat}
           \end{equation*}
        \partsitem           
           $\colvec{x \\ y}
           \mapsto
           \begin{mat}
             -3x/5-4y/5 \\
             -4x/5+3y/5
           \end{mat}
           +\colvec[r]{1 \\ 1}
           =\begin{mat}
             -3x/5-4y/5+1 \\
             -4x/5+3y/5+1
           \end{mat}$
      \end{exparts}
    \end{answer}
  \item \label{exer:IsometryFacts}
     \begin{exparts}
       \partsitem Pembuktian bahwa suatu pemetaan yang mempertahankan jarak dan
         membawa vektor nol ke dirinya sendiri sekaligus menunjukkan bahwa
         pemetaan tersebut injektif dan surjektif (titik dalam domain yang
         ditentukan oleh $d_0$, $d_1$, dan $d_2$ berkorespondensi dengan titik
         dalam kodomain yang ditentukan oleh ketiganya).
         Karena itu, sebarang pemetaan yang mempertahankan jarak memiliki
         invers. Tunjukkan bahwa invers tersebut juga mempertahankan jarak.
       \partsitem Buktikan bahwa kongruensi merupakan relasi ekuivalensi di
         antara bangun-bangun bidang.
     \end{exparts}
     \begin{answer}
       \begin{exparts}
         \partsitem Misalkan $f$ mempertahankan jarak dan tinjau $f^{-1}$.
           Sebarang dua titik dalam kodomain dapat ditulis sebagai $f(P_1)$
           dan $f(P_2)$. Karena $f$ mempertahankan jarak, jarak dari $f(P_1)$
           ke $f(P_2)$ sama dengan jarak dari $P_1$ ke $P_2$.
           Inilah tepatnya syarat agar $f^{-1}$ mempertahankan jarak.
         \partsitem Sebarang bangun bidang $F$ kongruen dengan dirinya sendiri
           melalui pemetaan identitas $\map{\identity}{\Re^2}{\Re^2}$, yang
           jelas mempertahankan jarak.
           Jika $F_1$ kongruen dengan $F_2$ (melalui suatu $f$), maka $F_2$
           kongruen dengan $F_1$ melalui $f^{-1}$, yang mempertahankan jarak
           menurut butir sebelumnya.
           Terakhir, jika $F_1$ kongruen dengan $F_2$ (melalui suatu $f$) dan
           $F_2$ kongruen dengan $F_3$ (melalui suatu $g$), maka $F_1$
           kongruen dengan $F_3$ melalui $\composed{g}{f}$, yang mudah
           diperiksa mempertahankan jarak.
       \end{exparts}
     \end{answer}
  \item \label{exer:HomoCrds}
    Dalam praktik, matriks bagi transformasi linear yang mempertahankan jarak
    dan translasinya sering digabungkan menjadi satu.
    Periksa bahwa kedua perhitungan berikut menghasilkan dua komponen pertama
    yang sama.
        \begin{equation*}
           \begin{mat}
             a  &c  \\
             b  &d
           \end{mat}
           \colvec{x  \\ y}
           +\colvec{e \\ f}
           \qquad
           \begin{mat}
             a  &c  &e \\
             b  &d  &f \\
             0  &0  &1
           \end{mat}
           \colvec{x  \\ y \\ 1}
        \end{equation*}
    (Ini disebut \definend{koordinat homogen}\index{koordinat homogen}; lihat
     Topik tentang Geometri Proyektif).
     \begin{answer}
       Dua komponen pertama masing-masing adalah $ax+cy+e$ dan $bx+dy+f$.
     \end{answer}
  \item \label{exer:GeomInvDistPre}
    \begin{exparts}
     \partsitem Verifikasi bahwa sifat-sifat yang disebut invarian terhadap
       pemetaan yang mempertahankan jarak pada paragraf kedua Topik ini memang
       bersifat demikian.
     \partsitem Berdasarkan pengalaman Anda mempelajari geometri Euklides,
       berikan dua sifat menarik lainnya yang juga invarian terhadap pemetaan
       yang mempertahankan jarak.
     \partsitem Berikan suatu sifat yang tidak menjadi perhatian dalam
       geometri Euklides dan tidak invarian terhadap pemetaan yang
       mempertahankan jarak.
    \end{exparts}
    \begin{answer}
      \begin{exparts}
        \partsitem Kasus kesamaan dalam ketaksamaan segitiga menyatakan bahwa
          tiga titik kolinear jika
          dan hanya jika (untuk suatu pengurutan menjadi $P_1$, $P_2$, dan
          $P_3$),
          $\dist(P_1,P_2)+\dist(P_2,P_3)=\dist(P_1,P_3)$.
          Tentu saja, jika $f$ mempertahankan jarak, hal ini berlaku jika dan
          hanya jika
          $\dist(f(P_1),f(P_2))+\dist(f(P_2),f(P_3))=\dist(f(P_1),f(P_3))$,
          yang, sekali lagi menurut kasus kesamaan ketaksamaan segitiga, benar
          jika dan hanya jika
          $f(P_1)$, $f(P_2)$, dan $f(P_3)$ kolinear.
        
          Argumen bagi keterantaraan serupa (di atas, $P_2$ berada di antara
          $P_1$ dan $P_3$).

          Jika bangun $F$ merupakan segitiga, maka $F$ adalah gabungan tiga
          ruas garis $P_1P_2$, $P_2P_3$, dan $P_1P_3$.
          Kedua paragraf sebelumnya bersama-sama menunjukkan bahwa sifat
          sebagai ruas garis bersifat invarian.
          Jadi, $f(F)$ merupakan gabungan tiga ruas garis dan karena itu
          merupakan segitiga.

          Lingkaran $C$ yang berpusat di $P$ dan berjari-jari $r$ merupakan
          himpunan semua titik $Q$ sedemikian sehingga $\dist(P,Q)=r$.
          Penerapan pemetaan $f$ yang mempertahankan jarak membuat bayangan
          $f(C)$ menjadi himpunan semua $f(Q)$ dengan syarat
          $\dist(P,Q)=r$. Karena $\dist(P,Q)=\dist(f(P),f(Q))$, himpunan
          $f(C)$ juga merupakan lingkaran, dengan pusat $f(P)$ dan jari-jari
          $r$.
        \partsitem Berikut dua sifat yang mudah diverifikasi: (i)~sifat sebagai
          segitiga siku-siku, dan (ii)~sifat dua garis sebagai garis sejajar.
        \partsitem Salah satu sifat yang disebut dalam bagian ini adalah
          `orientasi' suatu bangun.
          Segitiga yang titik-titik sudutnya terbaca searah jarum jam sebagai
          $P_1$, $P_2$, $P_3$ dapat dipetakan oleh pemetaan yang mempertahankan
          jarak ke segitiga yang $P_1$, $P_2$, $P_3$-nya terbaca berlawanan
          arah jarum jam.
      \end{exparts}
    \end{answer}
\end{exercises}
\index{matriks!ortonormal|)}
\index{matriks ortonormal|)}
\endinput
